\AtBeginSection[]{}
\section{Annexe}

\begin{frame}[plain]
  \footnotesize \tableofcontents[currentsection,sectionstyle=show/hide,hideothersubsections,subsubsectionstyle=shaded/shaded/shaded/hide]
  \addtocounter{framenumber}{-1}
\end{frame}

\subsection{Code}

\addtocontents{toc}{\catcode95=11}
\catcode95=11

\tiny

\subsubsection{caml_light.grammar}
\texttt{caml_light.grammar}
\inputminted[breaklines=true]{text}{../../caml_light.grammar}
\pagebreak

\foreach \name in {
    automaton.ml,
    caml_light.ml,
    grammar_grammar.ml,
    lexer.ml,
    parser.ml,
    rectify_helper.ml,
    rectify.ml,
    regex_grammar.ml,
    regex.ml,
    utils.ml,
    while_cli.ml,
    while.ml}{
    \subsubsection{\name}
    \texttt{\name}
    \inputminted[]{ocaml}{../../src/\name}
    \pagebreak
  }

\catcode95=8
\addtocontents{toc}{\catcode95=8}

\subsection{Algorithmes}

\subsubsection{Rectifier}

\tiny \begin{algorithm}[H]
  \caption{Rectifier}
  \Input{
    Une forme linéaire correcte $l$ rectifiée par $E$ ne contenant pas de
    variables nommées \texttt{cont} ou \texttt{new\_cont}.
  }
  \Output{
    Une forme linéaire correcte $l'$ rectifiable par $E$ telle que toute
    fonction de $E$ définie dans $l'$ soit récursive terminale, et pour
    toute fonction $f$ et en notant $a$ la dernière variable de
    $l$, $l.[(res, f\ a)]$ est équivalent à $[(\texttt{cont}, f)].l'$.
  }

  $(a, e) \leftarrow$ le premier élément de $l$ \\
  $q \leftarrow$ le reste de $l$ \\

  \uSi{$e$ contient une définition d'une fonction de $E$}{
    Mettre à jour $e$ tel que la fonction définie dedans:
    \begin{itemize}
      \item Prenne un argument \texttt{cont} en plus
      \item Évalue $Rectifier(l', E)$ lors de son appel o\`u $l'$
            est la forme linéaire évaluée \\ avant cette redéfinition.
    \end{itemize}
  }

  \bigskip

  \uSi{$e$ ne contient pas d'appel à une fonction de $E$}{
    \uSi{$q$ est vide}{
      Renvoyer
      \begin{flalign*}
        a   & := e                & \\
        res & := \texttt{cont } a &
      \end{flalign*}
    }\uSinon{
      $q_{rec} \leftarrow$ Rectifier($q$, $E$) \\
      Renvoyer $[(a, e)].q_{rec}$
    }
  }

  \ldots
\end{algorithm}

\tiny \begin{algorithm}[H]
  \ldots \\
  \uSinon{
    \uSi{$e$ est un appel de fonction}{
      $e_{rec} \leftarrow$ $e$ avec \texttt{cont} rajouté en dernier
      argument
    }\uSinon{
      $e_{rec} \leftarrow$ $e$ avec chaque enfant localement terminal $c$
      remplacé par Rectifier($c$, $E$)
    }

    \bigskip

    \uSi{$q$ est vide}{
      Renvoyer $[(a, e_{rec})]$
    }\uSinon{
      $q_{rec} \leftarrow$ Rectifier($q$, $E$) \\
      Renvoyer
      \begin{flalign*}
        \texttt{new\_cont} & := \texttt{fun } a \texttt{ -> } q_{rec} & \\
        \texttt{cont}      & := \texttt{new\_cont}                    & \\
        a                  & := e_{rec}                               &
      \end{flalign*}
    }
  }
\end{algorithm}

\pagebreak

\subsection{Plus d'exemples}

\begin{minted}{ocaml}
let rec sum (l : int list) : int = match l with [] -> 0 | x :: q -> x + sum q


let sum_rectified_whilified (l : int list) (acc : int) : int =
  let res_ref = ref None in
  let acc_ref = ref acc in
  let l_ref = ref l in
  while !res_ref = None do
    let acc = !acc_ref in
    let l = !l_ref in
    match l with
    | [] -> res_ref := Some acc
    | x :: q ->
        let new_acc (a_7 : int) : int = x + a_7 in
        l_ref := q;
        acc_ref := new_acc acc
  done;
  Option.get !res_ref

let sum (l : int list) : int = sum_rectified_whilified l 0
\end{minted}




\begin{minted}{ocaml}
type int_tree = Leaf of int | Node of int_tree * int * int_tree

let rec sum_tree (t : int_tree) : int =
  match t with
  | Leaf v -> v
  | Node (left, v, right) -> sum_tree left + v + sum_tree right


let rec sum_tree_rectified (t : int_tree) (cont : int -> 'ret) : 'ret =
  match t with
  | Leaf v -> cont v
  | Node (left, v, right) ->
      let new_cont (a_6 : int) : 'ret =
        let new_cont (a_15 : int) : 'ret = cont (a_6 + v + a_15) in
        sum_tree_rectified right new_cont
      in
      sum_tree_rectified left new_cont

let sum_tree (t : int_tree) : int = sum_tree_rectified t (fun res -> res)
\end{minted}


\begin{minted}{ocaml}
let rec fibonacci (n : int) : int =
  if n = 0 || n = 1 then 1 else fibonacci (n - 1) + fibonacci (n - 2)

    
let rec fibonacci_rectified (n : int) (cont : int -> 'ret) : 'ret =
  if n = 0 || n = 1 then cont 1
  else
    let new_cont (a_20 : int) : 'ret =
      let new_cont (a_29 : int) : 'ret = cont (a_20 + a_29) in
      fibonacci_rectified (n - 2) new_cont
    in
    fibonacci_rectified (n - 1) new_cont

let fibonacci (n : int) : int = fibonacci_rectified n (fun res -> res)    
\end{minted}

\begin{minted}{ocaml}
let rec custom_fold_left (fn : 'acc -> 'el -> 'ac) (init : 'acc) (l : 'el list)
    : 'ac =
  match l with [] -> init | x :: q -> custom_fold_left fn (fn init x) q


let custom_fold_left_rectified_whilified (fn : 'acc -> 'el -> 'ac) (init : 'acc)
    (l : 'el list) (cont : 'ac -> 'ret) : 'ret =
  let res_ref = ref None in
  let cont_ref = ref cont in
  let l_ref = ref l in
  let init_ref = ref init in
  let fn_ref = ref fn in
  while !res_ref = None do
    let cont = !cont_ref in
    let l = !l_ref in
    let init = !init_ref in
    let fn = !fn_ref in
    match l with
    | [] -> res_ref := Some (cont init)
    | x :: q ->
        fn_ref := fn;
        init_ref := fn init x;
        l_ref := q;
        cont_ref := cont
  done;
  Option.get !res_ref

let custom_fold_left (fn : 'acc -> 'el -> 'ac) (init : 'acc) (l : 'el list) :
    'ac =
  custom_fold_left_rectified_whilified fn init l (fun res -> res)
\end{minted}

\begin{minted}{ocaml}
let rec sum_even_numbers (l : int list) : int =
  match l with
  | [] -> 0
  | x :: q -> if x mod 2 = 0 then x + sum_even_numbers q else sum_even_numbers q


let sum_even_numbers_rectified_whilified (l : int list) (acc : int) : int =
  let res_ref = ref None in
  let acc_ref = ref acc in
  let l_ref = ref l in
  while !res_ref = None do
    let acc = !acc_ref in
    let l = !l_ref in
    match l with
    | [] -> res_ref := Some acc
    | x :: q ->
        if x mod 2 = 0 then (
          let new_acc (a_16 : int) : int = x + a_16 in
          l_ref := q;
          acc_ref := new_acc acc)
        else (
          l_ref := q;
          acc_ref := acc)
  done;
  Option.get !res_ref

let sum_even_numbers (l : int list) : int =
  sum_even_numbers_rectified_whilified l 0
\end{minted}

\begin{minted}{ocaml}
let rec is_even_length (l : 'a list) : bool =
  match l with [] -> true | _ :: q -> not (is_even_length q)


let is_even_length_rectified_whilified (l : 'a list) (acc : bool) : bool =
  let res_ref = ref None in
  let acc_ref = ref acc in
  let l_ref = ref l in
  while !res_ref = None do
    let acc = !acc_ref in
    let l = !l_ref in
    match l with
    | [] -> res_ref := Some acc
    | _ :: q ->
        let new_acc (a_6 : bool) : bool = not a_6 in
        l_ref := q;
        acc_ref := new_acc acc
  done;
  Option.get !res_ref

let is_even_length (l : 'a list) : bool =
  is_even_length_rectified_whilified l true
\end{minted}

\begin{minted}{ocaml}
let rec length (l : 'a list) : int =
  match l with [] -> 0 | _ :: q -> 1 + length q

let length_rectified_whilified (l : 'a list) (acc : int) : int =
  let res_ref = ref None in
  let acc_ref = ref acc in
  let l_ref = ref l in
  while !res_ref = None do
    let acc = !acc_ref in
    let l = !l_ref in
    match l with
    | [] -> res_ref := Some acc
    | _ :: q ->
        let new_acc (a_7 : int) : int = 1 + a_7 in
        l_ref := q;
        acc_ref := new_acc acc
  done;
  Option.get !res_ref

let length (l : 'a list) : int = length_rectified_whilified l 0
\end{minted}

\begin{minted}{ocaml}
let rec f () : unit =
  let rec g () : unit = f () in
  g ()


let rec f_rectified () (cont : unit -> 'ret) : 'ret =
  let rec g_rectified () (cont : unit -> 'ret) : 'ret = f_rectified () cont in
  g_rectified () cont

let f () : unit = f_rectified () (fun res -> res)
\end{minted}

\begin{minted}{ocaml}

type 'a tree = Leaf of 'a | Node of 'a tree list

let rec visit_tree (f : 'a -> unit) (t : 'a tree) : unit =
  let rec iter_list (g : 'b -> unit) (l : 'b list) : unit =
    match l with
    | [] -> ()
    | x :: q ->
        g x;
        iter_list g q
  in
  match t with
  | Leaf l -> f l
  | Node children -> iter_list (fun child -> visit_tree f child) children


let rec visit_tree_rectified (f : 'a -> unit) (t : 'a tree)
    (cont : unit -> 'ret) : 'ret =
  let rec iter_list_rectified (g_rectified : 'b -> (unit -> 'ret) -> 'ret)
      (l : 'b list) (cont : unit -> 'ret) : 'ret =
    match l with
    | [] -> cont ()
    | x :: q ->
        let new_cont (a_28 : unit) : 'ret =
          iter_list_rectified g_rectified q cont
        in
        g_rectified x new_cont
  in
  match t with
  | Leaf l -> cont (f l)
  | Node children ->
      let a_9_rectified child cont = visit_tree_rectified f child cont in
      iter_list_rectified a_9_rectified children cont

let visit_tree (f : 'a -> unit) (t : 'a tree) : unit =
  visit_tree_rectified f t (fun res -> res)
\end{minted}

\subsection{Fonctions récursives terminales}

\addtocounter{framenumber}{-1}
\begin{frame}[fragile,t]
  \tiny \begin{minted}{ocaml}
(* Non récursif terminal *)
let rec fibo (n : int) : int =
    if n = 0 || n = 1 then 1
    else
        (* L'addition est faite après le retour des appels récursifs *)
        fibo (n - 1) + fibo (n - 2)

(* Récursif terminal *)
let fibo (n : int) : int =
    let rec fibo_depuis (i : int) (f_i : int) (f_im1 : int) : int =
        if i = n then f_n
        else
            (* Cet appel récursif est la dernière opération effectuée *)
            fibo_depuis (i + 1) (f_i + f_im1) f_i
    in
    fibo_depuis 0 1 0
    \end{minted}
\end{frame}

\small

\subsection{Preuve de l'indécidabilité de ENS\_RECTIFIANT}
Par réduction depuis le co-problème de l'arrêt.

Soit $(M, w)$ une entrée de CO-ARRÊT. Considérons la forme linéaire suivante,
notée $l_{(M, w)}$:
\begin{align*}
  a_1 & := \text{Une sérialisation de $M$} \\
  a_2 & := \text{Une sérialisation de $w$} \\
  a_3 & := \texttt{m\_univ}                \\
  a_4 & := a_3\  a_1\  a_2                 \\
  a_5 & := \texttt{f}                      \\
  a_6 & := \texttt{()}                     \\
  a_7 & := a_5 \ a_6
\end{align*}

Avec \texttt{m\_univ} la machine universelle et \texttt{f} une fonction qui
conduit à un appel récursif de cette fonction.

Cette forme linéaire est rectifiable par $\emptyset$ si et seulement si $M$
ne s'arrête pas sur $w$.

On a alors
$(M, w) \in \text{CO-ARRÊT}^+ \iff (l_{(M, w)}, 0) \in \text{ENS\_RECTIFIANT}^+$.

Ainsi le problème CO-ARRÊT se réduit à ENS\_RECTIFIANT donc ENS\_RECTIFIANT
est indécidable.
